% ============================================================
% ch11：结语（对应论文 §8 Conclusion + 项目导航）
% ============================================================

\section{结语：这条研究路线的启示}\label{sec:ch11}

\begin{keybox}{本章速览}
\begin{itemize}
  \item 一句话总结：廉价尖锐度\emph{预测}成立（单 HVP 监测器 +
        梯度探针），廉价尖锐度\emph{控制}失败（帽的三种命运：原地、
        向上、有动作仍被钉住；唯全局动作下行）。
  \item 论文结尾金句：“The Edge of Stability supports a sharp
        diagnostic, but not a cheap steering wheel.”——EoS 支持一个
        锐利的诊断工具，但不支持一个廉价的方向盘。
  \item 三类读者各取所需：实践者用 $\lhatg$ 监测、优化器研究者避开
        曲率控制的单方向死路、理论工作者接手 G1--G5。
  \item 项目导航：FINDINGS\_zh.md（34 条发现的演化日志，含被自我推翻
        的中间结论）$\to$ tutorial/（入门）$\to$ paper/（本文）。
\end{itemize}
\end{keybox}

\textbf{一段式回顾。}这条研究路线从两边同时逼近同一个问题。预测端
（第~\ref{sec:ch08}~章之前的论文部分）：在 plain、BatchNorm、LayerNorm
架构上，廉价的尖锐度 proxy 可靠追踪 $\lmax$，其顶点是单 HVP 的梯度瑞利
商 $\lhatg$——在 EoS 对齐区近乎精确，且正确的探针是原始梯度
$\mathbf{g}$ 而非动量平滑的更新方向。控制端（第~\ref{sec:ch08}~
与~\ref{sec:ch09}~章）：自稳定帽给了 EoS 均衡一个明确的机构——周期-2
振荡本身——于是每种执行器干预的结局都可以\emph{预先}枚举：阻尼顶部
方向，帽向上移（上逃逸到 $261.8\pm5.4$，测试损失反而最好）；阻尼正交
补或用滞后的梯度方向传感器，帽原地不动（钉在 $2/\eta$）；在规模上
（ResNet-18/CIFAR-100），修正感知后的 rank-1 plant 全力开火却仍被钉住
——\emph{有动作仍被钉住}。三种命运，没有一种向下；向下只有一条路——
SAM 的全局曲率放大动作（三向对比 $22.63\pm1.73$ vs
$29.21\pm1.90/28.83\pm1.67$，CI 不重叠），而忠实的 X-SAM 消融
（$n=10$）把 SAM 的效力归到全局成分：其 rank-1 修正每步开火、效应量
上界 $+1.85$ 低于噪声地板、无精度增益。理论把这幅图像钉在 Damian 二维
递推上（Claim 1 + 命题 1：无翻号周期-2 轨道 $y^{*}\le0$），并以
G1--G5 诚实标出从二维约化定理到轨迹定理的距离。

\textbf{金句解读。}论文的最后一句话是：“The Edge of Stability supports
a sharp diagnostic, but not a cheap steering wheel.”——EoS 支持一个
锐利的诊断工具，但不支持一个廉价的方向盘。这句话的两个形容词都经过
校准：“锐利”，因为单次 HVP 就能在 EoS 体制里近乎精确地读出 $\lmax$
（去趋势相关 $>0.98$、对齐 $0.85$--$0.9$）；“廉价”不成立，因为把
$\lmax$ 挪出均衡的动作（SAM 一族）必须同时在\emph{每个}方向上施力，
其成本结构正是“不廉价”的。诊断与方向盘的落差不是工程欠账，而是
结构事实——观测器-控制器分离原理的类比之所以贴切，正在于“系统可以
容许便宜的观测器、却不容许同价位的控制器”。理解了帽的机构
（第~\ref{sec:ch09}~章），这句话就从警句变成了定理的推论。

\textbf{对不同读者的行动建议。}
\begin{itemize}
  \item \emph{实践者}：把 $\lhatg$（1 HVP）当作部署级的 $\lmax$
        监测器；在 EoS 对齐区（$|\cos(\mathbf{g},\vtope)|>0.7$，
        渐进尖锐化会自动把你带进去）信任它的相对走势，绝对值记得
        ${\sim}17\%$ 的下界松弛；对齐区之外用 $\|\mathbf{g}\|$ 兜底。
        不要部署任何“读个 proxy、阻尼一个方向”的曲率控制器。
  \item \emph{优化器研究者}：曲率感知的单方向干预这条路在秩 1 处关闭，
        且失败可由理论预言；剩下的路有两条——全局动作（SAM/ASAM/SAF 族）或
        步长调谐（CDAT、warmup 调度器）。同时记住局限 (vi) 的提醒：
        “更尖”未必“更差”，尖锐度与泛化的耦合是任务依赖的。
  \item \emph{理论工作者}：G1--G5（多特征值解耦、旋转奇点、时变
        $\gamma$、高阶 $|x|$-界、感知误差）是从二维约化定理到轨迹定理
        的现成问题清单；多尖锐景观上的逐方向独立是可形式化的实证
        起点；随机/动量体制下“阈值会移”的 EoS 理论同样有待建立。
\end{itemize}

\textbf{与教程收官的呼应。}教程的收官是“看懂 EoS”：$\lmax$ 爬向
$2/\eta$、周期-2 振荡、自稳定化的定性图像、对齐签名、轨迹 PCA、proxy
与普适性（universality，同一套定量规律在 MLP/CNN/Transformer 等不同
架构上反复出现的现象级一致性）。本文的收官是在同一片地基上竖起一块工程界碑：EoS \emph{能用
来看}，\emph{不能用来开}。从教程的现象观察到本文的工程边界，构成一条
完整的闭环——而这条闭环的特别之处在于，它是用一系列设计良好、诚实
报告的\emph{负结果}闭合的。

\textbf{项目导航。}若你想继续深入这个项目：\texttt{FINDINGS\_zh.md} 是
全部 34 条发现的带日期演化日志，包括那些\emph{被后续发现自我推翻}的
中间结论——X-SAM “0\% 触发”的假象、ResNet-18 幂迭代伪收敛造成的
“批尖锐度在 $2/\eta$”误读——它们是理解“负结果如何被正确地做出来”的
最佳教材；\texttt{tutorial/} 是 EoS 入门（本解读假设你已读完）；
\texttt{paper/} 是论文本身。三份材料对应同一条证据链的三个高度：
现象、方法、边界。
